% Rozdział aktywny po restrukturyzacji DT.
\section{Transformacja Fouriera w przestrzeniach Schwartza i funkcji całkowalnych}
\providecommand{\atlasref}[2]{#2}

\subsection{Transformacja Fouriera w przestrzeni Schwartza}

\noindent\emph{Powiązania tematyczne w Bibliotece matematycznej:} \atlasref{PDE-08::thm:pde8-inversion}{PDE-08: twierdzenie o odwracaniu transformacji Fouriera}.


W całym rozdziale $d\ge1$ jest liczbą całkowitą,
pracujemy nad $\C$ i przyjmujemy stałą
$c_d=(2\pi)^{-d/2}$. Piszemy
$\mathcal S=\mathcal S(\R^d;\C)$ i
$L^p=L^p(\R^d;\C)$, o ile nie wskazano inaczej.
Korzystamy z seminorm $s_{\alpha,\beta}$ i $Q_{N,m}$
przestrzeni Schwartza z rozdziału~\ref{sec:gladkie-schwartz17}.
Wzór z całką jest zwykłą całką Lebesgue'a,
bez potrzeby odwoływania się do dystrybucji.

\begin{definicja}[Transformacja Fouriera]
\label{def:fourier-schwartz17}
Dla $f\in\mathcal S(\R^d;\C)$ definiujemy
\[
\widehat f(\xi)=\mathcal Ff(\xi)=
(2\pi)^{-d/2}\int_{\R^d}e^{-ix\cdot\xi}f(x)\,
\mathrm d\lambda_d(x),\qquad\xi\in\R^d.
\]
Całka jest bezwzględnie zbieżna, gdyż $\mathcal S\subseteq L^1$
na mocy Stwierdzenia~\ref{prop:schwartz-operators17}.
\end{definicja}

\begin{twierdzenie}[Transformacja Fouriera zachowuje $\mathcal S$]
\label{thm:fourier-schwartz17}
Dla każdego $d\ge1$ transformacja Fouriera jest ciągłym
izomorfizmem liniowo-topologicznym
$\mathcal F:\mathcal S(\R^d;\C)\to\mathcal S(\R^d;\C)$.
Jej odwrotność wyraża wzór
\[
\mathcal F^{-1}g(x)=(2\pi)^{-d/2}
\int_{\R^d}e^{ix\cdot\xi}g(\xi)\,
\mathrm d\lambda_d(\xi).
\]
Dla każdego $f\in\mathcal S$ oraz $j=1,\dots,d$ zachodzą
$\mathcal F(\partial_jf)(\xi)=i\xi_j\widehat f(\xi)$
i $\mathcal F(x_jf)(\xi)=i\partial_{\xi_j}\widehat f(\xi)$.
\end{twierdzenie}
\begin{proof}
Przy różniczkowaniu względem $\xi_j$ iloraz różnicowy
czynnika wykładniczego jest ograniczony przez $|x_j|$.
Indukcyjnie majorantą przy kolejnych różniczkowaniach
jest $|x^{\beta+e_j}f(x)|$, całkowalna na mocy
Stwierdzenia~\ref{prop:schwartz-operators17}.
Twierdzenie o zbieżności zmajoryzowanej daje więc
$\partial_\xi^\beta\widehat f=(-i)^{|\beta|}
\mathcal F(x^\beta f)$, a także ciągłość tych pochodnych.
Dla pochodnej $\partial_j f$ stosujemy twierdzenie Fubiniego
i całkowanie przez części względem $x_j$ na $[-R,R]$.
Przy ustalonych pozostałych współrzędnych wyrazy brzegowe
dążą do zera, gdyż $f$ szybko maleje; całki jednowymiarowe
są bezwzględnie zbieżne. Następnie całkujemy po pozostałych
współrzędnych, korzystając z całkowalności $f$ i $\partial_j f$.
Otrzymujemy $\mathcal F(\partial_j f)=i\xi_j\widehat f$.
Iteracja daje $\mathcal F(\partial^\alpha f)=(i\xi)^\alpha\widehat f$.
Połączenie tych tożsamości prowadzi do
\[
\xi^\alpha\partial_\xi^\beta\widehat f(\xi)
=(-i)^{|\alpha|+|\beta|}
\mathcal F\bigl(\partial^\alpha(x^\beta f)\bigr)(\xi).
\]
Wartość bezwzględną prawej strony ograniczamy przez
$(2\pi)^{-d/2}\|\partial^\alpha(x^\beta f)\|_{L^1}$.
Z oszacowań w Stwierdzeniu~\ref{prop:schwartz-operators17}
otrzymujemy, dla pewnej stałej $C_{\alpha,\beta,d}$,
$s_{\alpha,\beta}(\widehat f)\le
C_{\alpha,\beta,d}Q_{d+1+|\beta|,|\alpha|}(f)$,
stosując kolejno oszacowanie normy $L^1$, różniczkowania
i mnożenia przez wielomian. Dokładniej,
\[
 \|\partial^\alpha(x^\beta f)\|_1
 \le C_d Q_{d+1,0}(\partial^\alpha(x^\beta f))
 \le C_d Q_{d+1,|\alpha|}(x^\beta f)
 \le C_{\alpha,\beta,d}'Q_{d+1+|\beta|,|\alpha|}(f).
\]
Tutaj $C_d=\int_{\R^d}(1+|x|)^{-d-1}\,\mathrm d\lambda_d(x)$
jest skończona na mocy tego samego stwierdzenia.
Wszystkie seminormy $s_{\alpha,\beta}(\widehat f)$ są więc
skończone i zależą w sposób ciągły od $f$.

Pozostaje udowodnić wzór odwrotny. Dla $\varepsilon>0$ definiujemy
\[
I_\varepsilon(x)=(2\pi)^{-d/2}\int_{\R^d}
 e^{-\varepsilon|\xi|^2/2}e^{ix\cdot\xi}
 \widehat f(\xi)\,\mathrm d\lambda_d(\xi).
\]
Ponieważ $f\in L^1$ i czynnik Gaussa należy do $L^1$,
możemy zastosować twierdzenie Fubiniego. Jednowymiarową całkę Gaussa obliczamy następująco:
$I(u)=\int_\R e^{-\varepsilon t^2/2}e^{iut}\,\mathrm d\lambda_1(t)$
można różniczkować pod całką z majorantą
$|t|e^{-\varepsilon t^2/2}$. Ponieważ
$(e^{-\varepsilon t^2/2})'=-\varepsilon t e^{-\varepsilon t^2/2}$,
całkowanie przez części, z zanikającymi wyrazami brzegowymi, daje
\[
I'(u)=-\frac{i}{\varepsilon}
\int_\R (e^{-\varepsilon t^2/2})'e^{iut}\,\mathrm d\lambda_1(t)
=-\frac u\varepsilon I(u).
\]
Ponadto $I(0)=\sqrt{2\pi/\varepsilon}$ z całki Gaussa
(po podstawieniu $s=\sqrt\varepsilon t$).
Dla przypomnienia, jeżeli
$J=\int_\R e^{-s^2/2}\,\mathrm d\lambda_1(s)$,
to $0<J<\infty$ (dla $|s|\ge1$ porównujemy funkcję
podcałkową z $e^{-|s|/2}$).
Z twierdzenia Tonellego i podstawienia biegunowego
\[
 J^2=\int_{\R^2} e^{-(s^2+t^2)/2}\,\mathrm d\lambda_2(s,t)
 =2\pi\int_0^\infty re^{-r^2/2}\,\mathrm d\lambda_1(r)
 =2\pi.
\]
Ostatnią całkę obliczamy podstawieniem $v=r^2/2$.
Z równania różniczkowego dla $I$ wynika
$(e^{u^2/(2\varepsilon)}I(u))'=0$, zatem
$I(u)=\sqrt{2\pi/\varepsilon}e^{-u^2/(2\varepsilon)}$.
Zastosowanie tego wzoru w każdej współrzędnej daje
\[
I_\varepsilon(x)=\int_{\R^d}f(y)
(2\pi\varepsilon)^{-d/2}e^{-|x-y|^2/(2\varepsilon)}
\,\mathrm d\lambda_d(y).
\]
Oznaczmy $G_\varepsilon(z)=(2\pi\varepsilon)^{-d/2}
e^{-|z|^2/(2\varepsilon)}$. Podstawienie $z=\sqrt\varepsilon w$
daje $\int G_\varepsilon\,\mathrm d\lambda_d=1$ oraz
$\int_{|z|\ge\delta}G_\varepsilon(z)\,\mathrm d\lambda_d(z)
=\int_{|w|\ge\delta/\sqrt\varepsilon}G_1(w)\,\mathrm d\lambda_d(w)
\to0$ dla każdego $\delta>0$, z całkowalności $G_1$.
Dla ustalonego $x$ i $\eta>0$ wybieramy $\delta>0$ tak,
aby $|f(x-z)-f(x)|<\eta$ dla $|z|<\delta$.
Wtedy
\[
|I_\varepsilon(x)-f(x)|\le
\eta+2\|f\|_\infty
\int_{|z|\ge\delta}G_\varepsilon(z)\,\mathrm d\lambda_d(z).
\]
Stąd granica górna błędu nie przekracza $\eta$, a ponieważ
$\eta$ jest dowolne, $I_\varepsilon(x)\to f(x)$.
Z drugiej strony $\widehat f\in L^1$, zatem
z twierdzenia o zbieżności zmajoryzowanej
$I_\varepsilon(x)$ dąży do całki ze wzoru odwrotnego.
Otrzymujemy $\mathcal F^{-1}\mathcal Ff=f$.
Ten sam argument, ze zmienionym znakiem wykładnika,
daje $\mathcal F\mathcal F^{-1}g=g$ dla $g\in\mathcal S$.
Ciągłość operatora odwrotnego wynika z identycznych
oszacowań seminorm po zmianie znaku wykładnika.
\end{proof}

\begin{wniosek}[Odbicie i iteracje transformacji Fouriera]
\label{cor:dt-fourier-iteration}
Niech $R:\mathcal S(\R^d;\C)\to\mathcal S(\R^d;\C)$
będzie operatorem odbicia $(Rf)(x)=f(-x)$.
Wtedy
\[
 \mathcal F^{-1}=R\mathcal F=\mathcal F R,\qquad
 \mathcal F^2=R,\qquad \mathcal F^4=I.
\]
Ponadto dla $f\in\mathcal S$ zachodzi
\[
 \widehat{\overline f}(\xi)=\overline{\widehat f(-\xi)}.
\]
\end{wniosek}
\begin{proof}
Odbicie jest ciągłym automorfizmem $\mathcal S$,
ponieważ zachowuje wszystkie seminormy $Q_{N,m}$
i spełnia $R^2=I$. W istocie
$\partial^\alpha(Rf)(x)=(-1)^{|\alpha|}(\partial^\alpha f)(-x)$
i $|-x|=|x|$.
Z wzoru odwrotnego w Twierdzeniu~\ref{thm:fourier-schwartz17}
otrzymujemy $\mathcal F^{-1}=R\mathcal F$.
Podstawienie $y=-x$ w całce Fourierowskiej daje
$\mathcal F R=R\mathcal F$, ponieważ miara Lebesgue'a
jest niezmiennicza względem odbicia. Zatem
$I=\mathcal F^{-1}\mathcal F=R\mathcal F^2$,
a mnożenie przez $R$ daje $\mathcal F^2=R$.
Stąd $\mathcal F^4=R^2=I$.
Sprzężenie zespolone przenosimy pod znak bezwzględnie
zbieżnej całki, co daje ostatnią tożsamość.
\end{proof}

\begin{przyklad}[Funkcja Gaussa i wizualizacja w Plotly]
\label{ex:fourier-gauss17}
Dla $a>0$ oraz $f_a(x)=e^{-ax^2/2}$, $x\in\R$,
stosując obliczoną w poprzednim dowodzie całkę $I(u)$
z $\varepsilon=a$ i $u=-\xi$, otrzymujemy
$\widehat f_a(\xi)=a^{-1/2}e^{-\xi^2/(2a)}$.
W szczególności zwiększenie skupienia funkcji $f_a$
odpowiada rozszerzeniu jej transformaty Fouriera.
Poniższy kod Python rysuje obie funkcje z użyciem biblioteki Plotly.
\begin{lstlisting}[language=Python,caption={Para funkcji Gaussa: funkcja i transformata Fouriera}]
import numpy as np
import plotly.graph_objects as go

a = 2.0
x = np.linspace(-5.0, 5.0, 1000)
f = np.exp(-a * x**2 / 2)
f_hat = a**(-0.5) * np.exp(-x**2 / (2 * a))

fig = go.Figure()
fig.add_scatter(x=x, y=f, mode="lines", name="f(x)")
fig.add_scatter(x=x, y=f_hat, mode="lines", name="Fourier(f)(x)")
fig.update_layout(title="Para Gaussa", xaxis_title="x / xi",
                  yaxis_title="wartosc")
fig.show()
\end{lstlisting}
\end{przyklad}

\begin{stwierdzenie}[Splot funkcji szybko malejących]
\label{prop:convolution-schwartz17}
Dla $f,g\in\mathcal S(\R^d;\C)$ definiujemy
$(f*g)(x)=\int_{\R^d}f(y)g(x-y)\,\mathrm d\lambda_d(y)$.
Wówczas $f*g\in\mathcal S$, odwzorowanie
$(f,g)\mapsto f*g$ jest dwuliniowe i łącznie ciągłe
na $\mathcal S\times\mathcal S$, a ponadto
$\mathcal F(f*g)=(2\pi)^{d/2}\widehat f\,\widehat g$.
\end{stwierdzenie}
\begin{proof}
Z $\mathcal S\subset L^1\cap L^\infty$ wynika
bezwzględna zbieżność splotu, gdyż całka modułu
nie przekracza $\|f\|_1\|g\|_\infty$.
Ilorazy różnicowe pochodnych $g$ względem $x_j$
są ograniczone przez $\|\partial^{\beta+e_j}g\|_\infty$,
więc po pomnożeniu przez $|f(y)|$ otrzymujemy majorantę
całkowalną względem $y$. Indukcja i twierdzenie o zbieżności
zmajoryzowanej dają $\partial^\beta(f*g)=f*(\partial^\beta g)$
i ciągłość wszystkich tych pochodnych.
Nierówność $1+|x|\le(1+|y|)(1+|x-y|)$ prowadzi do
\[
 Q_{N,m}(f*g)\le Q_{N,m}(g)
 \int_{\R^d}(1+|y|)^N|f(y)|\,\mathrm d\lambda_d(y)
 \le C_{d,N}Q_{N+d+1,0}(f)Q_{N,m}(g).
\]
To dowodzi jednocześnie szybkiego zanikania i ciągłości
splotu. Dla jasności, przy ustalonych $f_0,g_0\in\mathcal S$
rozkładamy różnicę
\[
 f*g-f_0*g_0=(f-f_0)*(g-g_0)
 +(f-f_0)*g_0+f_0*(g-g_0).
\]
Powyższe oszacowanie zastosowane do trzech składników
pokazuje, że dla każdego $N,m$ seminorma różnicy
dąży do zera przy $(f,g)\to(f_0,g_0)$ w topologii
produktowej $\mathcal S\times\mathcal S$.
Twierdzenie Tonellego i niezmienniczość miary Lebesgue'a
względem przesunięć dają
$\int_{\R^d}\int_{\R^d}|f(y)g(x-y)|\,\mathrm d\lambda_d(x)
\,\mathrm d\lambda_d(y)=\|f\|_1\|g\|_1<\infty$.
Możemy więc zastosować twierdzenie Fubiniego i podstawić $z=x-y$:
\[
\begin{aligned}
\mathcal F(f*g)(\xi)
&=(2\pi)^{-d/2}
\int_{\R^d} f(y)e^{-iy\cdot\xi}\,\mathrm d\lambda_d(y)
\int_{\R^d}g(z)e^{-iz\cdot\xi}\,\mathrm d\lambda_d(z)\\
&=(2\pi)^{d/2}\widehat f(\xi)\widehat g(\xi).
\end{aligned}
\]
\end{proof}

\begin{stwierdzenie}[Transformata iloczynu i splotu]
\label{prop:dt-fourier-product}
Dla $f,g\in\mathcal S(\R^d;\C)$ zachodzą
\[
 \mathcal F(fg)=c_d(\widehat f*\widehat g),\qquad
 \mathcal F^{-1}(fg)=c_d(\mathcal F^{-1}f*\mathcal F^{-1}g),
 \qquad c_d=(2\pi)^{-d/2}.
\]
W szczególności transformacja Fouriera zamienia
splot na mnożenie punktowe z czynnikiem $c_d^{-1}$,
a mnożenie punktowe na splot z czynnikiem $c_d$.
\end{stwierdzenie}
\begin{proof}
Reguła Leibniza daje
$Q_{N,m}(fg)\le 2^m Q_{N,m}(f)Q_{0,m}(g)$,
więc $fg\in\mathcal S$. Rzeczywiście, w każdym składniku
pochodnej iloczynu rzędu co najwyżej $m$ przypisujemy
całą wagę $(1+|x|)^N$ pochodnej $f$, a suma
współczynników dwumianowych dla wielowskaźnika
$\alpha$ wynosi $2^{|\alpha|}\le2^m$.
Wstawiamy wzór odwrotny dla $g$ z
Twierdzenia~\ref{thm:fourier-schwartz17}:
\[
\begin{aligned}
 \widehat{fg}(\xi)
 &=c_d^2\int_{\R^d}\int_{\R^d}
 f(x)\widehat g(\eta)e^{-ix\cdot(\xi-\eta)}
 \,\mathrm d\lambda_d(\eta)\,\mathrm d\lambda_d(x)\\
 &=c_d\int_{\R^d}\widehat f(\xi-\eta)\widehat g(\eta)
 \,\mathrm d\lambda_d(\eta).
\end{aligned}
\]
Zamiana kolejności całkowania jest dozwolona na mocy
twierdzenia Fubiniego, ponieważ całka modułu bez czynnika
$c_d^2$ wynosi $\|f\|_1\|\widehat g\|_1<\infty$.
Druga równość wynika z pierwszej i Wniosku~\ref{cor:dt-fourier-iteration}:
odbicie spełnia $R(u*v)=(Ru)*(Rv)$ przez podstawienie
$y=-z$ w całce definiującej splot.
Regułę transformowania splotu udowodniono w
Stwierdzeniu~\ref{prop:convolution-schwartz17}.
\end{proof}

\begin{przyklad}[Gauss wielowymiarowy i zmiana skali]
\label{ex:dt-fourier-gaussian-matrix}
Niech $A$ będzie rzeczywistą symetryczną macierzą
dodatnio określoną rozmiaru $d\times d$.
Dla $f_A(x)=\exp(-x^TAx/2)$ mamy
\[
 \widehat f_A(\xi)=(\det A)^{-1/2}
 \exp(-\xi^TA^{-1}\xi/2).
\]
Istotnie, z twierdzenia spektralnego dla macierzy
symetrycznych istnieje macierz ortogonalna $U$
oraz $a_1,\ldots,a_d>0$ takie, że
$A=U^T\operatorname{diag}(a_1,\ldots,a_d)U$.
W szczególności $x^TAx\ge(\min_j a_j)|x|^2$,
więc $f_A\in\mathcal S$: każda jej pochodna jest
iloczynem wielomianu i $f_A$, co wynika indukcyjnie
z reguły iloczynu.
Podstawienie $y=Ux$ zachowuje miarę Lebesgue'a,
a $x\cdot\xi=y\cdot(U\xi)$.
Twierdzenie Fubiniego pozwala rozdzielić całkę na
$d$ jednowymiarowych całek (ich moduły są całkowalne).
Z Przykładu~\ref{ex:fourier-gauss17} otrzymujemy
\[
 \widehat f_A(\xi)
 =\prod_{j=1}^d a_j^{-1/2}
 \exp\left(-\frac{(U\xi)_j^2}{2a_j}\right).
\]
Ponieważ $\det A=\prod_j a_j$ i
$A^{-1}=U^T\operatorname{diag}(a_1^{-1},\ldots,a_d^{-1})U$,
dostajemy tezę. W szczególności kierunki silniejszego
skupienia funkcji odpowiadają kierunkom słabszego
zanikania transformaty.
\end{przyklad}

\subsection{Transformacja całkowa w przestrzeni \texorpdfstring{$L^1$}{L1}}

Przyjmujemy w całym skrypcie normalizację unitarną z
Definicji~\ref{def:fourier-schwartz17}. Rozróżniamy
transformację (operator) i transformatę (jego wartość).
W $L^1$ całka określa funkcję w każdym punkcie, natomiast
w $L^2$ twierdzenie Plancherela określa w pierwszej
kolejności klasę funkcji równych prawie wszędzie.
Dla $f\in L^1(\R^d)$ definiujemy
\[
 \widehat f(\xi)=c_d\int_{\R^d}f(x)e^{-ix\cdot\xi}
 \,\mathrm d\lambda_d(x),\qquad c_d=(2\pi)^{-d/2}.
\]
Całka jest bezwzględnie zbieżna dla każdego $\xi$.
Jeżeli dwie reprezentantki klasy $f$ są równe prawie
wszędzie, to ich iloczyny z wykładnikiem również
są równe prawie wszędzie dla każdego ustalonego $\xi$.
Transformata jest więc dobrze określona na klasach
funkcji w $L^1$, a dla $f\in\mathcal S$ definicje
pokrywają się.
Intuicyjnie jest to pomiar amplitudy i fazy poszczególnych
częstotliwości; przy dużej częstotliwości wkłady z różnych
punktów mogą się znosić, nawet gdy moduł jądra jest równy
jeden. To zjawisko opisuje następne twierdzenie.

\begin{twierdzenie}[Lemat Riemanna--Lebesgue'a]
\label{thm:dt-fourier-l1-rl}
Transformacja całkowa jest ograniczonym operatorem
liniowym z $L^1(\R^d)$ do $C_0(\R^d)$, gdzie $C_0$
oznacza przestrzeń ciągłych funkcji zanikających w
nieskończoności, z normą supremum. Zachodzi
$\|\widehat f\|_\infty\le c_d\|f\|_1$.
\end{twierdzenie}
\begin{proof}
Liniowość i oszacowanie wynikają bezpośrednio z całki.
Ciągłość jest konsekwencją twierdzenia o zbieżności
zmajoryzowanej z majorantą $|f|$.
Aby wykazać zanikanie, dla $h\in\R^d$ oznaczmy
$\tau_h f(x)=f(x-h)$. Z niezmienniczości miary
Lebesgue'a względem translacji
$\widehat{\tau_h f}(\xi)=e^{-ih\cdot\xi}\widehat f(\xi)$.
Dla $\xi\ne0$ wybieramy $h=\pi\xi/|\xi|^2$.
Wtedy $e^{-ih\cdot\xi}=-1$ oraz
\[
 2|\widehat f(\xi)|
 =|\widehat{f-\tau_h f}(\xi)|
 \le c_d\|f-\tau_h f\|_1.
\]
Gdy $|\xi|\to\infty$, mamy $h\to0$, a prawa strona
dąży do zera na mocy ciągłości translacji w $L^1$,
wykazanej w dowodzie Twierdzenia~\ref{thm:dt-approx-lp}.
Por. także
\atlasref{ME-14::thm:translation-lp-14}{ciągłość translacji w przestrzeniach Lebesgue'a}.
\end{proof}

\begin{wniosek}[Jednostajna ciągłość transformaty]
\label{cor:dt-fourier-uniform-continuity}
Dla $f\in L^1(\R^d)$ i $h\in\R^d$ zachodzi
\[
 \sup_{\xi\in\R^d}|\widehat f(\xi+h)-\widehat f(\xi)|
 \le c_d\int_{\R^d}|f(x)|\,|e^{-ix\cdot h}-1|
 \,\mathrm d\lambda_d(x).
\]
W szczególności $\widehat f$ jest jednostajnie ciągła.
\end{wniosek}
\begin{proof}
Odejmujemy od siebie całki definiujące transformatę
w punktach $\xi+h$ i $\xi$, wyłączając wspólny
czynnik $e^{-ix\cdot\xi}$ o module $1$.
Nierówność trójkąta pod całką daje oszacowanie
niezależne od $\xi$. Przy $h\to0$ funkcja podcałkowa
dąży punktowo do zera, a jej moduł nie przekracza
$2|f(x)|$. Twierdzenie o zbieżności zmajoryzowanej
dowodzi jednostajnej ciągłości.
\end{proof}

\begin{stwierdzenie}[Reguły transformacji całkowej]
\label{prop:dt-fourier-l1-rules}
Dla $f,g\in L^1(\R^d)$, $h,\eta\in\R^d$ i odwracalnej
rzeczywistej macierzy $A$ zachodzą tożsamości
\[
\begin{aligned}
 \widehat{\tau_h f}(\xi)&=e^{-ih\cdot\xi}\widehat f(\xi),\\
 \mathcal F(e^{i\eta\cdot x}f(x))(\xi)&=\widehat f(\xi-\eta),\\
 \mathcal F(f(A\,\cdot))(\xi)&=
 |\det A|^{-1}\widehat f(A^{-T}\xi),\\
 \widehat{f*g}(\xi)&=c_d^{-1}\widehat f(\xi)\widehat g(\xi).
\end{aligned}
\]
Jeżeli ponadto $x^\beta f\in L^1$ dla wszystkich
wielowskaźników $|\beta|\le m$, to $\widehat f\in C^m$
oraz $\partial_\xi^\beta\widehat f=(-i)^{|\beta|}
\mathcal F(x^\beta f)$.
\end{stwierdzenie}
\begin{proof}
Pierwsza tożsamość wystąpiła w dowodzie
Twierdzenia~\ref{thm:dt-fourier-l1-rl}; druga wynika z
połączenia wykładników. W trzeciej stosujemy podstawienie
$y=Ax$, dla którego $x\cdot\xi=y\cdot A^{-T}\xi$.
W ostatniej całkowalność modułu funkcji podcałkowej
wynika z
\[
 \int_{\R^d}\int_{\R^d}|f(y)g(x-y)|
 \,\mathrm d\lambda_d(y)\,\mathrm d\lambda_d(x)
 =\|f\|_1\|g\|_1.
\]
Twierdzenie Tonellego daje również bezwzględną zbieżność
całki definiującej $(f*g)(x)$ dla prawie każdego $x$
i oszacowanie $\|f*g\|_1\le\|f\|_1\|g\|_1$.
Możemy używać borelowskich reprezentantek $f,g$;
wtedy funkcja $(x,y)\mapsto f(y)g(x-y)$ jest mierzalna.
Zmiana reprezentantki nie zmienia splotu prawie
wszędzie, ponieważ zmiana o funkcję zerową prawie
wszędzie ma, z tego samego oszacowania, splot o normie
$L^1$ równej zeru.
Twierdzenie Fubiniego i podstawienie $z=x-y$ pozwalają
rozłożyć całkę na iloczyn dwóch całek. Czynnik wynosi
$c_d/c_d^2=c_d^{-1}$.
Dla pochodnych różniczkujemy pod całką, korzystając z
$|(e^{-itx_j}-1)/t|\le |x_j|$. Na każdym kolejnym etapie
majorantą jest $|x^{\beta+e_j}f(x)|$; ciągłość
pochodnych wynika ponownie ze zbieżności zmajoryzowanej.
\end{proof}

\begin{wniosek}[Transformata pochodnej słabej w $L^1$]
\label{cor:dt-fourier-weak-derivative-l1}
Niech $f,g\in L^1(\R^d;\C)$ i niech
$\partial_jT_f=T_g$ w $\mathcal D'(\R^d)$
dla ustalonego $j\in\{1,\ldots,d\}$.
Wtedy dla każdego $\xi\in\R^d$
\[
 \widehat g(\xi)=i\xi_j\widehat f(\xi).
\]
W szczególności jeśli wszystkie pochodne słabe
$\partial^\alpha f$ dla $|\alpha|\le m$ należą do $L^1$,
to $\mathcal F(\partial^\alpha f)=(i\xi)^\alpha\widehat f$
dla każdego takiego wielowskaźnika.
\end{wniosek}
\begin{proof}
Nie można bezpośrednio użyć funkcji wykładniczej
jako testu w definicji pochodnej dystrybucyjnej,
gdyż nie ma zwartego nośnika.
Z Lematu~\ref{lem:cutoffs17} wybieramy
$\chi\in\mathcal D(\R^d;\R)$, $0\le\chi\le1$,
równą $1$ w otoczeniu zera. Dla $R>0$ i
$\xi\in\R^d$ bierzemy test
$\varphi_R(x)=\chi(x/R)e^{-ix\cdot\xi}$.
Z definicji pochodnej dystrybucyjnej
\[
\begin{aligned}
 \int_{\R^d}g(x)\varphi_R(x)\,\mathrm d\lambda_d(x)
 &=i\xi_j\int_{\R^d} f(x)\varphi_R(x)\,\mathrm d\lambda_d(x)\\
 &\quad-\frac1R\int_{\R^d}f(x)(\partial_j\chi)(x/R)
 e^{-ix\cdot\xi}\,\mathrm d\lambda_d(x).
\end{aligned}
\]
Ostatni składnik ma moduł nie większy niż
$R^{-1}\|f\|_1\|\partial_j\chi\|_\infty$, więc zanika.
W dwóch pozostałych całkach stosujemy twierdzenie
o zbieżności zmajoryzowanej, ponieważ
$\chi(x/R)\to1$ punktowo, a majoranty to odpowiednio
$|g|$ i $|f|$.
Mnożąc otrzymaną równość przez $c_d$, dostajemy tezę.
Iterowanie tego wyniku po kolejnych pochodnych
pierwszego rzędu daje wzór dla wielowskaźników:
jeśli $|\alpha|>0$ i $\alpha_j>0$, stosujemy wzór
do pochodnej słabej $\partial^{\alpha-e_j}f$,
której pochodna w kierunku $j$ jest równa
$\partial^\alpha f$ jako dystrybucja regularna.
Obie funkcje należą do $L^1$ z założenia.
Indukcja rozpoczyna się od $\alpha=0$.
W szczególności z Twierdzenia~\ref{thm:dt-fourier-l1-rl}
mamy $\xi^\alpha\widehat f(\xi)\to0$ przy $|\xi|\to\infty$.
\end{proof}

\begin{lemat}[Odwracanie z regularyzacją Gaussa]
\label{lem:dt-fourier-gaussian-l1}
Niech $f\in L^1(\R^d)$ i $\varepsilon>0$.
Definiujemy jądro Gaussa
\[
 G_\varepsilon(x)=(2\pi\varepsilon)^{-d/2}
 e^{-|x|^2/(2\varepsilon)}.
\]
Wtedy dla każdego $x$ zachodzi
\[
 c_d\int_{\R^d}e^{ix\cdot\xi}
 e^{-\varepsilon|\xi|^2/2}\widehat f(\xi)
 \,\mathrm d\lambda_d(\xi)=(G_\varepsilon*f)(x).
\]
Ponadto $G_\varepsilon*f\to f$ w $L^1$ dla
$\varepsilon\downarrow0$.
\end{lemat}
\begin{proof}
Lewa strona jest bezwzględnie zbieżna na mocy
Twierdzenia~\ref{thm:dt-fourier-l1-rl}.
Podstawiamy definicję $\widehat f$ i stosujemy Fubiniego,
gdyż $|f(y)|e^{-\varepsilon|\xi|^2/2}$ jest całkowalne
na przestrzeni produktowej. Całkę po $\xi$ obliczamy
wzorem Gaussa z dowodu Twierdzenia~\ref{thm:fourier-schwartz17};
łącznie z czynnikiem $c_d^2$ daje to $G_\varepsilon(x-y)$.
Splot po prawej stronie jest określony dla każdego $x$,
ponieważ $G_\varepsilon$ jest ograniczona i $f\in L^1$.
Z tego samego dowodu wiemy, że $\int G_\varepsilon=1$
oraz całka $G_\varepsilon$ poza dowolną kulą o dodatnim
promieniu dąży do zera. Po podstawieniu $z=x-y$
możemy więc zapisać
\[
 (G_\varepsilon*f)(x)-f(x)=
 \int_{\R^d}G_\varepsilon(y)(f(x-y)-f(x))\,\mathrm d\lambda_d(y)
\]
dla prawie każdego $x$.
Nierówność trójkąta pod całką i twierdzenie Tonellego
dla nieujemnej funkcji
$G_\varepsilon(y)|f(x-y)-f(x)|$ dają
\[
 \|G_\varepsilon*f-f\|_1
 \le \int_{\R^d}G_\varepsilon(y)
 \|\tau_yf-f\|_1\,\mathrm d\lambda_d(y).
\]
Dla dowolnego $\delta>0$ część po $|y|<\delta$ nie
przekracza $\sup_{|y|<\delta}\|\tau_yf-f\|_1$,
a pozostała nie przekracza
$2\|f\|_1\int_{|y|\ge\delta}G_\varepsilon\,\mathrm d\lambda_d$.
Ciągłość translacji z dowodu Twierdzenia~\ref{thm:dt-approx-lp}
pozwala najpierw wybrać $\delta$,
a następnie zmniejszyć $\varepsilon$, tak aby obie
części były dowolnie małe.
\end{proof}

\begin{wniosek}[Punktowa zbieżność regularyzacji Gaussa]
\label{cor:dt-gaussian-pointwise}
Niech $f\in L^1(\R^d)$ i niech wybrana reprezentantka
$f$ będzie ciągła w punkcie $x_0\in\R^d$.
Wtedy
\[
 \lim_{\varepsilon\downarrow0}(G_\varepsilon*f)(x_0)
 =f(x_0),
\]
gdzie $G_\varepsilon$ jest jądrem z
Lematu~\ref{lem:dt-fourier-gaussian-l1}.
Nie zakładamy, że $\widehat f\in L^1$ ani że $f$
jest ograniczona na całej przestrzeni.
\end{wniosek}
\begin{proof}
Dla $\eta>0$ wybieramy $\delta>0$ tak, aby
$|f(x_0-y)-f(x_0)|<\eta$ dla $|y|<\delta$.
Korzystamy z $\int G_\varepsilon\,\mathrm d\lambda_d=1$.
Nierówność trójkąta i rozbicie całki dają
\[
 \begin{aligned}
 |(G_\varepsilon*f)(x_0)-f(x_0)|
 &\le \eta+\|f\|_1
 \sup_{|y|\ge\delta}G_\varepsilon(y)\\
 &\quad+|f(x_0)|\int_{|y|\ge\delta}G_\varepsilon(y)
 \,\mathrm d\lambda_d(y).
 \end{aligned}
\]
W drugim składniku użyliśmy niezmienniczości miary
Lebesgue'a względem translacji i odbicia, aby ograniczyć
całkę z $|f(x_0-y)|$ przez $\|f\|_1$.
Dla ustalonego $\delta>0$
\[
 \sup_{|y|\ge\delta}G_\varepsilon(y)
 =(2\pi\varepsilon)^{-d/2}e^{-\delta^2/(2\varepsilon)}
 \longrightarrow0,
\]
gdyż funkcja wykładnicza zanika szybciej niż dowolna
ujemna potęga $\varepsilon$.
Trzeci składnik dąży do zera na mocy własności
jądra Gaussa wykazanej w dowodzie
Lematu~\ref{lem:dt-fourier-gaussian-l1}.
Granica górna błędu nie przekracza więc $\eta$,
a $\eta>0$ jest dowolna.
\end{proof}

\begin{twierdzenie}[Odwrócenie i jednoznaczność w $L^1$]
\label{thm:dt-fourier-l1-inversion}
Jeżeli $f,\widehat f\in L^1(\R^d)$, to
\[
 f(x)=c_d\int_{\R^d}e^{ix\cdot\xi}\widehat f(\xi)
 \,\mathrm d\lambda_d(\xi)
\]
dla prawie każdego $x$, a prawa strona jest ciągłym,
zanikającym w nieskończoności reprezentantem $f$.
Jeżeli $f$ jest ciągła, równość zachodzi wszędzie.
Niezależnie od całkowalności $\widehat f$ transformacja
na $L^1$ jest różnowartościowa.
\end{twierdzenie}
\begin{proof}
Niech $g$ oznacza prawą stronę wzoru odwrotnego.
Twierdzenie~\ref{thm:dt-fourier-l1-rl}, ze zmienionym
znakiem argumentu, daje $g\in C_0$.
Z Lematu~\ref{lem:dt-fourier-gaussian-l1} oraz
zbieżności zmajoryzowanej mamy
\[
 \|G_\varepsilon*f-g\|_\infty
 \le c_d\int_{\R^d}
 (1-e^{-\varepsilon|\xi|^2/2})|\widehat f(\xi)|
 \,\mathrm d\lambda_d(\xi)\longrightarrow0.
\]
Z drugiej strony $G_\varepsilon*f\to f$ w $L^1$.
Na każdej kuli skończonej miary zbieżność jednostajna
implikuje zbieżność w $L^1$ na tej kuli. Jednoznaczność
granicy daje $g=f$ prawie wszędzie na każdej takiej kuli,
a ich przeliczalne wyczerpanie daje równość na $\R^d$.
Jeżeli $f$ jest ciągła, różnica dwóch funkcji ciągłych
równa zeru prawie wszędzie jest identycznie zerowa.
Jeżeli $\widehat f=0$, nawet bez założenia o
całkowalności transformaty, lewa strona tożsamości
z Lematu~\ref{lem:dt-fourier-gaussian-l1} jest zerowa
dla każdego $\varepsilon$, więc $G_\varepsilon*f=0$.
Przejście do granicy w $L^1$ daje $f=0$.
Liniowość implikuje różnowartościowość.
\end{proof}

\begin{twierdzenie}[Obraz transformacji całkowej w $C_0$]
\label{thm:dt-fourier-l1-range}
Dla każdego $d\ge1$ operator
$\mathcal F:L^1(\R^d)\to C_0(\R^d)$ ma normę $c_d$.
Jego obraz jest gęstą, właściwą podprzestrzenią liniową
$C_0(\R^d)$, a operator odwrotny określony na tym obrazie
nie jest ciągły względem normy supremum.
\end{twierdzenie}
\begin{proof}
Oszacowanie $\|\mathcal F\|\le c_d$ pochodzi z
Twierdzenia~\ref{thm:dt-fourier-l1-rl}. Dla nieujemnej
funkcji $f\in L^1$ o normie $1$ mamy
$\widehat f(0)=c_d$, więc zachodzi równość.

Z Twierdzenia~\ref{thm:fourier-schwartz17}
obraz zawiera $\mathcal S$, w szczególności
$C_c^\infty(\R^d)$.
Aby wykazać gęstość, niech $v\in C_0(\R^d)$ i $\eta>0$.
Wybieramy $R>0$ tak, aby $|v(x)|<\eta$ poza $B(0,R)$,
oraz, z Lematu~\ref{lem:cutoffs17}, funkcję
$\chi\in\mathcal D(\R^d;\R)$ równą $1$ na
$\overline{B(0,R)}$ i spełniającą $0\le\chi\le1$.
Funkcja $w=\chi v$ jest ciągła o zwartym nośniku
i $\|w-v\|_\infty\le\eta$.
Jest ona jednostajnie ciągła: na dostatecznie dużej
domkniętej kuli wynika to ze zwartości, a dla dwóch
punktów o odległości mniejszej niż $1$, z których
jeden leży poza kulą zawierającą jednostkowe otoczenie
nośnika, obie wartości są zerowe.
Niech $\rho_\varepsilon$ będzie aproksymacją jedności
z Definicji~\ref{def:mollifier21}.
Z Twierdzenia~\ref{thm:smoothing21} splot
$w*\rho_\varepsilon$ jest funkcją testową, a
\[
 \|w*\rho_\varepsilon-w\|_\infty
 \le\sup_{|y|\le\varepsilon}\sup_x|w(x-y)-w(x)|
 \longrightarrow0.
\]
Dowodzi to gęstości.

Dla wykazania braku ciągłości odwrotności skonstruujemy
ciąg funkcji o normie $L^1$ równej $1$, których transformaty
dążą jednostajnie do zera.
Wybieramy nieujemne $\psi\in C_c^\infty(\R^d)$ o normie
$L^1$ równej $1$ oraz o nośniku w $B(0,1)$.
Dla $N=2^k$, $k\ge0$, określamy wielomiany rekurencyjnie:
\[
 P_0(z)=Q_0(z)=1,\qquad
 \begin{cases}
 P_{k+1}(z)=P_k(z)+z^{2^k}Q_k(z),\\
 Q_{k+1}(z)=P_k(z)-z^{2^k}Q_k(z).
 \end{cases}
\]
Indukcyjnie oba wielomiany mają stopień $2^k-1$
i wszystkie współczynniki równe $1$ lub $-1$:
przesunięcie o $2^k$ rozdziela zakresy potęg obu składników.
Dla $|z|=1$ tożsamość
$|a+b|^2+|a-b|^2=2(|a|^2+|b|^2)$ daje
\[
 |P_k(z)|^2+|Q_k(z)|^2=2^{k+1}=2N.
\]
Jeżeli $P_k(z)=\sum_{j=0}^{N-1}b_jz^j$,
połóżmy
\[
 f_k(x)=N^{-1}\sum_{j=0}^{N-1}b_j\psi(x-3je_1).
\]
Nośniki składników są parami rozłączne, stąd
$\|f_k\|_1=N^{-1}\sum_j|b_j|\|\psi\|_1=1$.
Z reguły translacji w
Stwierdzeniu~\ref{prop:dt-fourier-l1-rules}
\[
 \widehat f_k(\xi)=N^{-1}\widehat\psi(\xi)
 P_k(e^{-3i\xi_1}),
 \qquad
 \|\widehat f_k\|_\infty\le c_d\sqrt{2/N}\longrightarrow0.
\]
Różnowartościowość z Twierdzenia~\ref{thm:dt-fourier-l1-inversion}
pozwala mówić o odwrotności na obrazie. Powyższy ciąg
dowodzi, że nie jest ona ciągła w zerze.
Gdyby obraz był całym $C_0(\R^d)$, twierdzenie
o odwzorowaniu otwartym dla przestrzeni Banacha
$L^1$ i $C_0$ wymuszałoby ciągłość odwrotności.
Sprzeczność dowodzi, że obraz jest właściwy.
\end{proof}

\begin{przyklad}[Prostokąt i funkcja wykładnicza]
\label{ex:dt-fourier-l1-pairs}
Dla $a>0$ w wymiarze jeden mamy
\[
 \mathcal F\chi_{[-a,a]}(\xi)
 =\sqrt{\frac{2}{\pi}}\frac{\sin(a\xi)}{\xi},
 \qquad
 \mathcal F(e^{-a|x|})(\xi)
 =\sqrt{\frac{2}{\pi}}\frac{a}{a^2+\xi^2}.
\]
Pierwszą wartość przy $\xi=0$ rozumiemy jako
$\sqrt{2/\pi}\,a$. Pierwszy wzór wynika z całkowania
$e^{-ix\xi}$ po $[-a,a]$. W drugim rozdzielamy całkę
na półproste, otrzymując
$c_1((a+i\xi)^{-1}+(a-i\xi)^{-1})$.
Transformata prostokąta nie należy do $L^1$: na
przedziałach $a\xi\in[n\pi+\pi/6,n\pi+5\pi/6]$
dla $n\ge1$ licznik ma moduł co najmniej $1/2$,
więc suma całek jest ograniczona z dołu przez
dodatnią stałą razy $\sum_{n\ge1}(n+1)^{-1}$.
W drugim przypadku obie funkcje należą do $L^1$,
a wzór odwrotny obowiązuje w każdym punkcie.
\end{przyklad}

\begin{uwaga}[Zgodność z normalizacją RRC]
\label{rem:dt-fourier-normalization-rrc}
W skrypcie RRC-08 stosowana jest transformata
$\mathcal F_0f(\xi)=\int e^{-ix\cdot\xi}f(x)
\,\mathrm d\lambda_d(x)$. Mamy $\mathcal F=c_d\mathcal F_0$;
we wzorze odwrotnym dla $\mathcal F_0$ występuje
$(2\pi)^{-d}$, a we wzorze na splot nie ma dodatkowego
czynnika. To nie są odmienne twierdzenia, lecz różne
normalizacje tego samego operatora. Por.
\atlasref{PDE-08::prop:pde8-RL}{RRC-08: lemat Riemanna--Lebesgue'a}
oraz
\atlasref{PDE-08::thm:pde8-inversion}{RRC-08: odwrócenie transformaty}.
\end{uwaga}
